\chapter{Introdução}
\label{cap:intro}



A ascensão da Indústria 4.0 e da \gls{iot} consolidou a integração entre \emph{hardware} e \emph{software} como um pilar fundamental da automação moderna. Nesse cenário, a robótica móvel surge como uma das vertentes mais dinâmicas da engenharia, unificando mecânica, eletrônica e computação para a criação de sistemas capazes de interagir e deslocar-se de forma coordenada em ambientes físicos reais.

No âmbito acadêmico e de formação tecnológica, a utilização de plataformas robóticas móveis consolidou-se como uma ferramenta pedagógica de elevado valor experimental. De acordo com \citeonline{pinto2025juca}, o contato prático com esses sistemas permite aos estudantes de engenharia a aplicação direta de conceitos complexos de sistemas embarcados, teoria de controle e cinemática. Robôs dotados de rodas apresentam uma complexidade mecânica otimizada quando comparados a sistemas baseados em pernas (\emph{legged robots}), tornando-se os modelos ideais para laboratórios de ensino devido à sua alta manobrabilidade, replicabilidade de ensaios e facilidade de isolamento de variáveis dinâmicas~\cite{siegwart2011introduction}.

Posto isto, a plataforma robótica Juca foi desenvolvida no âmbito do \gls{daeln} do \gls{ifsc} como uma ferramenta de ensino e pesquisa voltada à robótica móvel~\cite{pinto2025juca}. Por se tratar de um ecossistema educacional concebido em ambiente acadêmico, sua consolidação e usabilidade prática demandam a contínua evolução de suas ferramentas de interatividade e supervisão. Nesse contexto, este trabalho tem como propósito expandir a operabilidade da plataforma Juca por meio do desenvolvimento de um aplicativo móvel dedicado, capaz de estabelecer a comunicação sem fio para a transmissão de comandos de controle e a recepção contínua de telemetria em tempo real.



\section{Justificativa}

A robótica educacional propõe o uso de robôs como recurso didático que possibilita aos alunos se envolverem em temas que estão inclusive fora do campo da robótica (CRISTOFORIS et al., 2013).
O robô móvel Juca é uma plataforma didática relevante para o estudo prático de sistemas embarcados e engenharia de controle. Entretanto, a utilização plena e eficiente dessa plataforma em ensaios acadêmicos exige uma camada de interação com o usuário que seja intuitiva e capaz de exibir variáveis de telemetria em tempo real.
Dessa forma, este trabalho justifica-se pelo projeto e desenvolvimento de um sistema para comunicação e  focado na criação de um aplicativo móvel interativo para o controle e monitoramento do robô móvel Juca

\section{Definição do Problema}

Plataformas robóticas didáticas, como o robô móvel Juca, necessitam de meios eficientes para o envio de comandos de movimentação e a recepção de dados de telemetria em tempo real durante ensaios práticos. No entanto, a ausência de uma interface de controle móvel dedicada e intuitiva restringe a operação do robô a conexões cabeadas ou a interfaces genéricas e pouco flexíveis. Essa limitação compromete a usabilidade do equipamento em ambientes de teste e dificulta a visualização dinâmica de variáveis operacionais essenciais, como velocidade, estado dos sensores.
Além disso, o estabelecimento de uma comunicação sem fio entre dispositivos móveis e sistemas embarcados impõe desafios técnicos em relação à latência de resposta, ao consumo energético e à estabilidade da transmissão de dados. 

\section{Objetivo Geral}

Projetar e implementar uma interface móvel interativa para o controle operacional e monitoramento de telemetria em tempo real da plataforma robótica móvel Juca.

\section{Objetivos Específicos}

\begin{alineas}
    \item Projetar a interface gráfica do usuário do aplicativo móvel, desenvolvendo componentes interativos para o envio de comandos de movimentação e visualização dos dados do robô Juca.
    \item Implementar a arquitetura de comunicação sem fio no microcontrolador ESP32-S3 da plataforma embarcada, configurando o protocolo Bluetooth Low Energy (BLE) para o envio e recebimento de pacotes de dados. 
    \item Desenvolver a lógica de integração e o tratamento de dados no aplicativo para estabelecer  conexão, leitura e escrita de características BLE com baixa latência e tratamento de desconexões.
    \item Realizar ensaios práticos e testes de validação com o robô móvel Juca, avaliando o desempenho do sistema integrado em termos de estabilidade da conexão, latência de resposta dos comandos e precisão do monitoramento.
\end{alineas}
    
\section{Estrutura do Trabalho}

Este trabalho está organizado da seguinte forma: o Capítulo 1 apresenta a contextualização do tema, a justificativa do projeto, a definição do problema, os objetivos e a estrutura do documento. O Capítulo 2 reúne a fundamentação teórica necessária para a compreensão dos conceitos de robótica móvel, a arquitetura do microcontrolador ESP32-S3, o protocolo Bluetooth Low Energy (BLE) e o desenvolvimento de interfaces reativas em Flutter. O Capítulo 3 descreve a metodologia adotada, incluindo o levantamento de requisitos do sistema, a concepção da interface do aplicativo, a implementação do protocolo BLE na plataforma embarcada e a lógica do joystick no aplicativo. O Capítulo 4 apresenta o desenvolvimento do ecossistema de software, detalhando a arquitetura de comunicação proposta, o fluxo de envio e recepção de pacotes e a integração entre o aplicativo móvel e o robô Juca. O Capítulo 5 apresenta os resultados experimentais de bancada e ensaios práticos, com as análises de latência, estabilidade do link sem fio e desempenho da telemetria em tempo real. Por fim, o Capítulo 6 reúne as conclusões do trabalho e as sugestões para pesquisas futuras.
